%-------------------------------------------------------------------------------
%	SECTION TITLE
%-------------------------------------------------------------------------------
\cvsection{Research}


%-------------------------------------------------------------------------------
%	CONTENT
%-------------------------------------------------------------------------------
\begin{cventries}

%---------------------------------------------------------
  \cventry
    {Post-doctoral Researcher} % Job title
    {\href{https://hci.isir.upmc.fr/}{CNRS} - \href{https://ismm.ircam.fr/}{IRCAM} - \href{https://ex-situ.lri.fr/}{INRIA}} % Organization
    {Paris, France} % Location
    {Jan. 2020 - Sept. 2022} % Date(s)
    {
      \begin{cvitems} % Description(s) of tasks/responsibilities
        \item {Developed a functionning prototype (Javascript, MaxMSP) for movement exploration and modeling, demoed at \textit{\href{https://nime.pubpub.org/pub/8zdsd4ar/}{NIME'22}}.}
        \item {Studied the perception of motor learning (Pytorch) in tasks involving complex and creative motions, published in \textit{\href{https://hal.archives-ouvertes.fr/hal-03790829/}{PLOS ONE}}.}
        \item {Used reinforcement learning (MABs) for optimising a training schedule in a motor task, submitted to IHM'22.}
        % https://journals.plos.org/plosone/article?id=10.1371/journal.pone.0272509
      \end{cvitems}
    }

%---------------------------------------------------------
  \cventry
    {Research Assistant and PhD Candidate} % Job title
    {\href{https://www.gla.ac.uk/schools/computing/research/researchsections/ida-section/inferencedynamicsandinteraction/}{University of Glasgow}} % Organization
    {Glasgow, United Kingdom} % Location
    {Mar. 2015 - Sep. 2019} % Date(s)
    {
      \begin{cvitems} % Description(s) of tasks/responsibilities
        % \item {Assessment of undergraduate projects and tutoring courses. Correction of examination essays.}
        % \item {Member of the Inference, Dynamics and Interaction (IDI) group led by Pr. Murray-Smith.}
        \item {Researched computational methods for the study of novel gestural interactions (\textit{\href{https://theses.gla.ac.uk/78981/}{thesis}}). Use of discriminative models to create virtual surfaces from RGB-D sensors (\textit{\href{https://dl.acm.org/doi/abs/10.1145/3132272.3135074}{ISS'17}}); optimisation of gameplay experiences for custom game controllers via user probabilistic modelling - on \textit{\href{https://arxiv.org/abs/2209.14788}{arXiv}}; audio-based reinforcement learning for the exploration and characterisation of user physical capabilities.}
        \item {Contributed to the European project MoreGrasp: lead redaction of deliverable documents, implementation of visual feedback mechanisms on smartwatches (Android), preparation of live demos and posters for review meetings, establishment of contacts with health practitioners at Glasgow's hospital leading to workshops and a funding proposal.}
        \item {Analysed visualisation techniques (\textit{\href{https://github.com/toinsson/raviz}{Jupyter notebooks}}) for low-dimensional embeddings (t-SNE).}
        \item {Tutored the course Artificial Intelligence and Data Fundamentals for Undergraduate and Master students.}
        % \item {reference: Pr. Chalmers - matthew.chalmers@glasgow.ac.uk}
      \end{cvitems}
    }

%---------------------------------------------------------
  \cventry
    {Research Assistant} % Job title
    {Stockholm University, \textit{\href{http://www.mobilelifecentre.org/}{Mobile Life}}} % Organization
    {Stockholm, Sweden} % Location
    {Feb. 2014 - Oct. 2014} % Date(s)
    {
      \begin{cvitems} % Description(s) of tasks/responsibilities
        % \item {Member of Pr. Brown’s group dedicated to the understanding of human interactions with novel technology through the use of ethnographic methods.}
        \item {Explored innovative designs for always-on in the background speech interactions, published at \textit{\href{https://dl.acm.org/doi/abs/10.1145/2702123.2702532}{CHI'15}}. Designed an experiment involving day-long recordings of participants’ everyday life with qualitative analysis and design workshop involving trained Interaction designers.}
        % \item {reference: Pr. Brown - barry@dsv.su.se}
      \end{cvitems}
    }
%---------------------------------------------------------
\end{cventries}



%---------------------------------------------------------
  % \cventry
  %   {PhD Candidate and Contributor to the European Project MoreGrasp} % Job title
  %   {University of Glasgow} % Organization
  %   {Glasgow, United Kingdom} % Location
  %   {Mar. 2015 - Feb. 2019} % Date(s)
  %   {
  %     \begin{cvitems} % Description(s) of tasks/responsibilities
  %       \item {Member of the Inference, Dynamics and Interaction (IDI) group led by Pr. Murray-Smith.}
  %       \item {Research involving computational methods for the study of novel gestural interactions: use of discriminative models to create virtual surfaces from RGB-D sensors, optimisation of gameplay experiences for custom game controllers via user probabilistic modelling, audio-based reinforcement learning for the exploration and characterisation of user physical capabilities. Partly published at \textit{\href{https://dl.acm.org/doi/abs/10.1145/3132272.3135074}{ISS’17}}.}
  %       \item {Contributions to the European project MoreGrasp: lead redaction of deliverable documents, implementation of visual feedback mechanisms on smartwatches, preparation of live demos and posters for review meetings, establishment of contacts with practitioners in the QEU hospital in Glasgow leading to workshops in-situ and the writing of a funding proposal targeting a local charity.}
  %       \item {Tutoring for Artificial Intelligence and Data Fundamentals courses (undergraduate, Master levels).}
  %       % \item {references: Pr. Murray-Smith - rod@glasgow.ac.uk, Dr. Williamson - jhw@glasgow.ac.uk}
  %     \end{cvitems}
  %   }